\chapter{Finite Automata and Regular Languages}

\section{Finite Automata}

\begin{greybox}
\begin{definition}[Deterministic Finite Automaton]
A \emph{deterministic finite automaton} (DFA) is a $5$-tuple $(Q, \Sigma, \delta, q_0, F)$, where
\begin{enumerate}
    \item $Q$ is a finite set called \emph{states},
    \item $\Sigma$ is a finite set called the \emph{alphabet},
    \item $\delta : Q \times \Sigma \to Q$ is a \emph{transition function},
    \item $q_0 \in Q$ is the \emph{start state}, and
    \item $F \subseteq Q$ is the \emph{set of accept states}.
\end{enumerate}
\end{definition}
\end{greybox}

\begin{example}
Look at the state diagram of the following finite automaton $M_1$.

\begin{center}
\begin{tikzpicture}[>=stealth, node distance=3cm, thick]
    \node[state, initial] (q0) {$q_0$};
    \node[state, accepting, right of=q0] (q1) {$q_1$};
    \node[state, right of=q1] (q2) {$q_2$};

    \draw (q0) edge[loop above] node{0} (q0)
          (q0) edge[above] node{1} (q1)
          (q1) edge[loop above] node{1} (q1)
          (q1) edge[bend left, above] node{0} (q2)
          (q2) edge[bend left, below] node{0} (q1)
          (q2) edge[loop above] node{1} (q2);
\end{tikzpicture}
\end{center}

Now a finite automaton takes a binary string, such as $1101$, as input and gives out the output if the input is \underline{accept} or \underline{reject}.

\medskip
\noindent\textbf{Q.} How to quantify this accept/reject?

Reading through the string from left to right, taking in each binary digit, the FA changes its state. If the last state is an accept state, such as $q_1$, then the output is accept, otherwise reject.

\bigskip
\noindent For this example, on input $1101$:

\begin{center}
\begin{tabular}{c c}
digit & state transition \\ \hline
$1$ & $q_0 \to q_1$ \\
$1$ & $q_1 \to q_1$ \\
$0$ & $q_1 \to q_2$ \\
$1$ & $q_2 \to q_1$
\end{tabular}
\end{center}

\noindent $\therefore\ 1101$ is accepted.
\end{example}

\bigskip
\noindent For this example,
\begin{enumerate}[(i)]
    \item $Q = \{q_0, q_1, q_2\}$
    \item $\Sigma = \{0, 1\}$
    \item $\delta : Q \times \Sigma \to Q$ with the following table
\end{enumerate}

\begin{center}
\begin{tabular}{c|c|c}
$\delta$ & $0$ & $1$ \\ \hline
$q_0$ & $q_0$ & $q_1$ \\
$q_1$ & $q_2$ & $q_1$ \\
$q_2$ & $q_1$ & $q_2$
\end{tabular}
\end{center}

\begin{enumerate}[(iv)]
    \item $q_0$ is the start state.
    \item $F = \{q_1\}$.
\end{enumerate}

\section{Languages}

The set of all strings $A$ that $M$ accepts is called the \emph{language of machine $M$}, denoted as $L(M)$. We say $M$ \emph{recognizes} $A = L(M)$.

A machine may accept many strings, but it only accepts one language. We also have the empty language ($\emptyset$) for machines accepting no strings.

\begin{example}
For $M_1$ as before, we have
\[
L(M_1) = \{ w \in \{0,1\}^* \mid w \text{ has at least one } 1 \text{ and an even number of } 0\text{'s follow the last } 1 \}.
\]
\end{example}

\begin{example}
Here is an example of the state diagram of a finite automaton $M_2, M_3, M_4$.

\medskip
\noindent\underline{$M_2$:}
\begin{center}
\begin{tikzpicture}[>=stealth, node distance=3cm, thick]
    \node[state, initial] (q0) {$q_0$};
    \node[state, accepting, right of=q0] (q1) {$q_1$};

    \draw (q0) edge[loop above] node{0} (q0)
          (q0) edge[bend left, above] node{1} (q1)
          (q1) edge[loop above] node{1} (q1)
          (q1) edge[bend left, below] node{0} (q0);
\end{tikzpicture}
\end{center}

\noindent Formally, $M_2 = (\{q_0, q_1\}, \{0,1\}, \delta, q_0, \{q_1\})$.

\bigskip
\noindent For this example, $\delta$ is given by:
\begin{center}
\begin{tabular}{c|c|c}
$\delta$ & $0$ & $1$ \\ \hline
$q_0$ & $q_0$ & $q_1$ \\
$q_1$ & $q_0$ & $q_1$
\end{tabular}
\end{center}

\noindent Now, language of $M_2$:
\[
L(M_2) = \{ w \mid w \text{ ends with } 1 \}.
\]

\bigskip
\noindent\underline{$M_3$:}
\begin{center}
\begin{tikzpicture}[>=stealth, node distance=3cm, thick]
    \node[state, initial, accepting] (q0) {$q_0$};
    \node[state, right of=q0] (q1) {$q_1$};

    \draw (q0) edge[loop above] node{0} (q0)
          (q0) edge[bend left, above] node{1} (q1)
          (q1) edge[loop above] node{1} (q1)
          (q1) edge[bend left, below] node{0} (q0);
\end{tikzpicture}
\end{center}

\[
L(M_3) = \{ w \mid w \text{ ends with } 0 \text{ or } w = \varepsilon \}.
\]

\bigskip
\noindent\underline{$M_4$:}
\begin{center}
\begin{tikzpicture}[>=stealth, node distance=2.6cm, thick]
    \node[state] (s) {$S$};
    \node[state, accepting, below left of=s] (q1) {$q_1$};
    \node[state, accepting, below right of=s] (r1) {$r_1$};
    \node[state, below of=q1] (q2) {$q_2$};
    \node[state, below of=r1] (r2) {$r_2$};

    \draw (s) edge[left] node{$a$} (q1)
          (s) edge[right] node{$b$} (r1)
          (q1) edge[loop left] node{$a$} (q1)
          (r1) edge[loop right] node{$b$} (r1)
          (q1) edge[bend left, left] node{$b$} (q2)
          (q2) edge[bend left, right] node{$a$} (q1)
          (q2) edge[loop left] node{$b$} (q2)
          (r1) edge[bend left, right] node{$a$} (r2)
          (r2) edge[bend left, left] node{$b$} (r1)
          (r2) edge[loop right] node{$a$} (r2);
\end{tikzpicture}
\end{center}

\[
Q = \{S, q_1, q_2, r_1, r_2\}, \qquad \Sigma = \{a, b\}, \qquad F = \{q_1, r_1\}, \qquad \text{start-state} = S.
\]

\noindent $L(M_4) = ?$

\medskip
It is easy to see that going from the left component to the right is not possible. And for the left (resp.\ right) component, the string must start with $a$ (resp.\ $b$), and for the string to be accepted, it must end with $a$ (resp.\ $b$). Also, we can change to $q_1$ (resp.\ $r_1$) iff the last string letter is $a$ (resp.\ $b$). So
\[
L(M_4) = \{ w \in \{a,b\}^* \mid w \text{ starts and ends with some (the same) letter} \}.
\]
\end{example}

\begin{greybox}
\begin{remark}
\begin{enumerate}[(1)]
    \item For a language, there may be different DFAs accepting it (in fact, there can be infinitely many). But each DFA accepts exactly one language.
    \item Given a pair $(\overset{\text{state}}{q}, \overset{\text{symbol}}{a})$, $\exists! \ q' \in Q$ so that $\delta(q,a) = q'$.
    \item The set of states in a DFA partitions the set of all strings. The set of states is finite.
\end{enumerate}
\end{remark}
\end{greybox}

\subsection{Computation of a DFA}

Let $M = (Q, \Sigma, \delta, q_0, F)$ and $w = a_1 \cdots a_n$ where $a_i \in \Sigma$. We say that $M$ \emph{accepts} $w$ if $\exists$ a sequence of states $r_0, \dots, r_n$ s.t.
\begin{enumerate}[1)]
    \item $r_0 = q_0$
    \item $r_i = \delta(r_{i-1}, a_i)\ \ \forall i = 1, \dots, n$
    \item $r_n \in F$.
\end{enumerate}

\begin{greybox}
\begin{remark}
The states $\{r_i\}$ need not be distinct.
\end{remark}
\end{greybox}

Let $L \subseteq \Sigma^*$ be a language. We say that $M$ \emph{accepts/recognizes} $L$ if
\[
L = \{ w \in \Sigma^* \mid M \text{ accepts } w \}.
\]

Let $L \subseteq \Sigma^*$ be a language. We say $L$ is \emph{regular} if $\exists$ a DFA $M$ s.t.
\[
L = L(M).
\]

\begin{example}
$A = \{10, 001\}$
\[
A^* = \{ \varepsilon, 10, 001, 10001, 1010, 00110, \dots \}.
\]
\end{example}

\section{Non-Determinism}

\begin{itemize}
    \item From a state $q$, on an input symbol $a$, the automaton can go to multiple states ($k$ different states, $k \geq 0$).
    \item Computation happens simultaneously along each of these paths.
    \item Has $\varepsilon$-transitions if there is an $\varepsilon$-transition
    \begin{center}
        \begin{tikzpicture}[>=stealth, node distance=3cm, thick]
            \node[state] (qi) {$q_i$};
            \node[state, right of=qi] (qj) {$q_j$};
            \draw (qi) edge[above] node{$\varepsilon$} (qj);
        \end{tikzpicture}
    \end{center}
    then the automaton moves to state $q_j$ from $q_i$ without reading the next input bit.
    \item An input is accepted if there is some computation path that leads to an accept state.
\end{itemize}

\begin{example}
\begin{center}
\begin{tikzpicture}[>=stealth, node distance=3cm, thick]
    \node[state, initial] (q0) {$q_0$};
    \node[state, right of=q0] (q1) {$q_1$};
    \node[state, right of=q1] (q2) {$q_2$};
    \node[state, accepting, right of=q2] (q3) {$q_3$};

    \draw (q0) edge[loop above] node{$0,1$} (q0)
          (q0) edge[above] node{$1$} (q1)
          (q1) edge[above] node{$0, \varepsilon$} (q2)
          (q2) edge[above] node{$1$} (q3)
          (q3) edge[loop above] node{$0,1$} (q3);
\end{tikzpicture}
\end{center}

\begin{align*}
(q_0, 1) &\to \{q_0, q_1\} \\
(q_1, 1) &\to \{\} \\
(q_1, 0) &\to \{q_2\} \\
(q_1, \varepsilon) &\to \{q_2\}
\end{align*}

\noindent\textbf{Note:} The transitions are labelled with symbols from the set $\Sigma_\varepsilon = \Sigma \cup \{\varepsilon\}$.
\end{example}

\bigskip
\begin{example}
$L = \{0^n 1^n \mid n \geq 0\}$ is not regular.
\end{example}

\begin{example}
$L = \{ w \mid w \text{ doesn't contain } 11 \text{ as a substring} \}$.

\begin{center}
\begin{tikzpicture}[>=stealth, node distance=3cm, thick]
    \node[state, initial, accepting] (q0) {$q_0$};
    \node[state, accepting, right of=q0] (q1) {$q_1$};
    \node[state, right of=q1] (q2) {$q_2$};

    \draw (q0) edge[loop above] node{0} (q0)
          (q0) edge[bend left, above] node{1} (q1)
          (q1) edge[bend left, below] node{0} (q0)
          (q1) edge[above] node{1} (q2)
          (q2) edge[loop above] node{$0,1$} (q2);
    \node[below=0.2cm of q2] {(dump state)};
\end{tikzpicture}
\end{center}

$q_2$ corresponds to the set of all strings that contain $11$ as a substring. $q_0$ corr.\ to the set of strings that don't end with a $1$. $q_1$ corr.\ to the set of all strings that don't contain $11$ and end with $1$.
\[
M = (\{q_0, q_1, q_2\}, \{0,1\}, \delta, q_0, \{q_0, q_1\}).
\]

\begin{greybox}
\begin{remark}
A \emph{dump state} is a state from where the automaton cannot reach an accept state.
\end{remark}
\end{greybox}
\end{example}

\begin{example}
$L = \{ w \in \{0,1\}^* \mid |w| = 3 \}$.

\begin{center}
\begin{tikzpicture}[>=stealth, node distance=2.6cm, thick]
    \node[state, initial] (q0) {$q_0$};
    \node[state, right of=q0] (q1) {$q_1$};
    \node[state, right of=q1] (q2) {$q_2$};
    \node[state, accepting, right of=q2] (q3) {$q_3$};
    \node[state, below of=q3] (q4) {$q_4$};

    \draw (q0) edge[above] node{$0,1$} (q1)
          (q1) edge[above] node{$0,1$} (q2)
          (q2) edge[above] node{$0,1$} (q3)
          (q3) edge[right] node{$0,1$} (q4)
          (q4) edge[loop below] node{$0,1$} (q4);
    \node[right=0.1cm of q4] {(dump state)};
\end{tikzpicture}
\end{center}
\end{example}

\subsection{Regular Operations}

Let $A, B$ be two languages over $\Sigma$.
\begin{enumerate}[1.]
    \item \textbf{Union:} $A \cup B = \{x \mid x \in A \text{ or } x \in B\}$
    \item \textbf{Concatenation:} $A \cdot B = \{xy \mid x \in A,\ y \in B\}$
    \item \textbf{Star:} $A^* = \{x_1 \cdots x_k \mid k \geq 0,\ x_i \in A\ \forall i\}$
\end{enumerate}

\section{Nondeterministic Finite Automata}

\begin{greybox}
\begin{definition}[NFA]
A \emph{nondeterministic finite automaton} (NFA) is the five-tuple $N = (Q, \Sigma, \delta, q_0, F)$ with $Q, \Sigma, q_0, F$ defined the same as for a DFA, but $\delta$ is different: instead of $\delta : Q \times \Sigma \to Q$, we have the transition function
\[
\delta : Q \times \Sigma_\varepsilon \to 2^Q, \qquad [\Sigma_\varepsilon = \Sigma \cup \{\varepsilon\}].
\]
\end{definition}
\end{greybox}

\begin{greybox}
\begin{definition}
We say that $N$ \emph{accepts} an input $w = a_1 \cdots a_n$ if we can write $w$ as $w = b_1 \cdots b_m$ where $b_i \in \Sigma_\varepsilon$ and $\exists$ a sequence of states $r_0, \dots, r_m$ (not necessarily distinct) such that
\begin{enumerate}[i)]
    \item $r_0 = q_0$ \hfill (initial condition)
    \item $r_i \in \delta(r_{i-1}, b_i) \ \ \forall i = 1, \dots, m$ \hfill (transition condition)
    \item $r_m \in F$. \hfill (accept condition)
\end{enumerate}
\end{definition}
\end{greybox}

We similarly define $L(N) := \{ w \in \Sigma^* \mid N \text{ accepts } w \}$.

\begin{example}
$L_3 = \{ w \in \{0,1\}^* \mid 3^{\text{rd}} \text{ last digit of } w \text{ is } 1 \}$.

\begin{center}
\begin{tikzpicture}[>=stealth, node distance=3cm, thick]
    \node[state, initial] (q0) {$q_0$};
    \node[state, right of=q0] (q1) {$q_1$};
    \node[state, accepting, right of=q1] (q2) {$q_2$};

    \draw (q0) edge[loop above] node{$0,1$} (q0)
          (q0) edge[above] node{1} (q1)
          (q1) edge[above] node{$0,1$} (q2);
\end{tikzpicture}
\end{center}

This choice (whether to stay at $q_0$ or move to $q_1$) is made by (nondeterministically) guessing if the digit read is the $3^{\text{rd}}$ last digit of $w$.

\bigskip
Similarly we can define an NFA for
\[
L_k = \{ w \in \{0,1\}^* \mid k^{\text{th}} \text{ last digit of } w \text{ is } 1 \}.
\]
\end{example}

\bigskip
\noindent Consider an input $w = 010110$.

\begin{center}
\begin{tabular}{c|c}
Set of states the automaton is in & Symbol read \\ \hline
$\{q_0\}$ & $0$ \\
$\{q_0\}$ & $1$ \\
$\{q_0, q_1, q_2\}$ & $0$ \\
$\{q_0, q_2\}$ & $1$ \\
$\{q_0, q_1, q_2, q_3\}$ & $1$ \\
$\{q_0, q_1, q_2, q_3\}$ & $0$ \\
$\{q_0, q_2, q_3\}$ & Input accepted
\end{tabular}
\end{center}

\subsection{Difference between DFA and NFA}

\begin{center}
\begin{tabular}{p{6.5cm} p{6.5cm}}
\textbf{DFA} & \textbf{NFA} \\ \hline
$(q,a) \to$ single state & $(q,a) \to$ multiple states \\
single computation & multiple computation \\
no $\varepsilon$-transitions & have $\varepsilon$-transitions \\
accept if the computation ends at an accept state & accept if one of the computation paths ends at an accept state
\end{tabular}
\end{center}

\section{Equivalence of NFA and DFA}

Every DFA is also an NFA.

\begin{blackbox}
\begin{theorem}
Let $N = (Q, \Sigma, \delta, q_0, F)$ be an NFA. There exists a DFA $D = (Q', \Sigma, \delta', q_0', F')$ such that $L(N) = L(D)$.
\end{theorem}
\end{blackbox}

\begin{proof}[Proof (construction of $D$)]
\begin{itemize}
    \item $Q' = 2^Q$.
    \item Let $A \subseteq Q$, then $\delta'(A, a) = \bigcup_{r \in A} \delta(r, a)$.
    \item $q_0' = \{q_0\}$.
    \item $F' = \{ A \subseteq Q \mid A \cap F \neq \emptyset \}$.
\end{itemize}

\noindent What about $\varepsilon$-transitions?

Let $R \subseteq Q$. Define the $\varepsilon$-\emph{closure} of $R$ as $E(R)$:
\[
E(R) = \bigcup_{r \in R} \{ t \mid t \text{ is reachable from } r \text{ using } 0 \text{ or more } \varepsilon\text{-transitions} \}.
\]

We define $\delta'$ to accommodate $\varepsilon$-transitions:
\[
\delta'(A, a) = \bigcup_{r \in A} E(\delta(r, a)).
\]

\noindent (Proof continued — refer to Hopcroft / Sipser / Kozen.)
\end{proof}

\begin{example}
$L_2 = \{ w \in \{0,1\}^* \mid |w| \text{ is div.\ by } 2 \text{ or } 5 \}$.
\end{example}

\noindent\textit{Solution.}
\begin{center}
\begin{tikzpicture}[>=stealth, node distance=2.2cm, thick]
    \node[state] (q0) {$q_0$};
    \node[state, accepting, above right of=q0] (q1) {$q_1'$};
    \node[state, right of=q1] (q2) {$q_2'$};
    \node[state, accepting, below right of=q0] (r1) {$q_1''$};
    \node[state, right of=r1] (r2) {$q_2''$};
    \node[state, right of=r2] (r3) {$q_3''$};
    \node[state, right of=r3] (r4) {$q_4''$};
    \node[state, below of=r3] (r5) {$q_5''$};

    \draw (q0) edge[above] node{$\varepsilon$} (q1)
          (q1) edge[bend left, above] node{0} (q2)
          (q2) edge[bend left, below] node{0} (q1)
          (q0) edge[below] node{$\varepsilon$} (r1)
          (r1) edge[above] node{0} (r2)
          (r2) edge[above] node{0} (r3)
          (r3) edge[above] node{0} (r4)
          (r4) edge[bend right, right] node{0} (r5)
          (r5) edge[bend right, left] node{0} (r1);
\end{tikzpicture}
\end{center}

\noindent $q_0$ moves via $\varepsilon$ into a $2$-cycle (tracking length mod $2$, accepting at the start of the cycle) and, separately, via $\varepsilon$ into a $5$-cycle (tracking length mod $5$, accepting at the start of that cycle).

\bigskip
\begin{example}
$L = \{ w \mid 2^{\text{nd}} \text{ last symbol in } w \text{ is } 1 \}$.

\begin{center}
\begin{tikzpicture}[>=stealth, node distance=3cm, thick]
    \node[state, initial] (q0) {$q_0$};
    \node[state, right of=q0] (q1) {$q_1$};
    \node[state, accepting, right of=q1] (q2) {$q_2$};

    \draw (q0) edge[loop above] node{$0,1$} (q0)
          (q0) edge[above] node{1} (q1)
          (q1) edge[above] node{$0,1$} (q2);
\end{tikzpicture}
\end{center}

\noindent\textbf{Construction of DFA for this} (subset construction):

\begin{center}
\begin{tikzpicture}[>=stealth, node distance=3.2cm, thick]
    \node[state, initial] (q0) {$q_0$};
    \node[state, right of=q0] (q01) {$q_{01}$};
    \node[state, accepting, below of=q01] (q012) {$q_{012}$};
    \node[state, accepting, right of=q012] (q02) {$q_{02}$};

    \draw (q0) edge[loop above] node{0} (q0)
          (q0) edge[bend left, above] node{1} (q01)
          (q01) edge[bend left, below] node{0} (q0)
          (q01) edge[right] node{0} (q02)
          (q01) edge[left] node{1} (q012)
          (q012) edge[loop below] node{$0,1$} (q012)
          (q02) edge[bend left, above] node{1} (q01)
          (q02) edge[bend right=60, below] node{0} (q0);
\end{tikzpicture}
\end{center}

\[
\delta(q_0,0) = q_0,\quad \delta(q_0,1) = q_{01}, \quad
\delta(q_{01},0) = q_{02}, \quad \delta(q_{01},1) = q_{012},
\]
\[
\delta(q_{02},0) = q_0, \quad \delta(q_{02},1) = q_{01}, \quad
\delta(q_{012},0) = q_{012}, \quad \delta(q_{012},1) = q_{012},
\]
\[
F = \{q_{02}, q_{012}\}.
\]

\begin{greybox}
\begin{remark}
States that are not reachable from the start state are known as \emph{dead states}.
\end{remark}
\end{greybox}
\end{example}

\section{Closure under Union, Concatenation, and Star}

\begin{blackbox}
\begin{theorem}
If $L_1$ and $L_2$ are regular, then $L_1 \cup L_2$, $L_1 L_2$, and $L_1^*$ are also regular.
\end{theorem}
\end{blackbox}

\begin{greybox}
\begin{remark}
For a regular language, we can assume that there is an NFA accepting it with a unique accept state (just use $\varepsilon$-transitions).
\end{remark}
\end{greybox}

\begin{proof}
Let $N_1, N_2$ be NFAs. Assume the following setup:

\begin{center}
\begin{tabular}{c|c|c|c}
Language & NFA & start state & end state \\ \hline
$L_1$ & $N_1$ & $q_1$ & $r_1$ \\
$L_2$ & $N_2$ & $q_2$ & $r_2$
\end{tabular}
\end{center}

\noindent with $L_1 = L(N_1)$ and $L_2 = L(N_2)$.

\medskip
\noindent\textbf{For $L_1 \cup L_2$:}
\begin{center}
\begin{tikzpicture}[>=stealth, thick, node distance=2.6cm]
    \node[state] (q) {};
    \node[draw, rectangle, minimum width=2.6cm, minimum height=1cm, above right of=q, xshift=1.2cm, yshift=0.3cm] (N1) {};
    \node at ($(N1.west)+(0.5,0)$) {$q_1$};
    \node at ($(N1.east)+(-0.5,0)$) {$r_1$};
    \node[draw, rectangle, minimum width=2.6cm, minimum height=1cm, below right of=q, xshift=1.2cm, yshift=-0.3cm] (N2) {};
    \node at ($(N2.west)+(0.5,0)$) {$q_2$};
    \node at ($(N2.east)+(-0.5,0)$) {$r_2$};
    \node[state, right of=N1, xshift=1.6cm, yshift=-1.1cm] (r) {};

    \draw (q) edge[above] node{$\varepsilon$} (N1.west)
          (q) edge[below] node{$\varepsilon$} (N2.west)
          (N1.east) edge[above] node{$\varepsilon$} (r)
          (N2.east) edge[below] node{$\varepsilon$} (r);
\end{tikzpicture}
\end{center}

So every string that is accepted by $L_1$ or $L_2$ is accepted for this setup as well.

\medskip
\noindent\textbf{For $L_1 L_2$:}
\begin{center}
\begin{tikzpicture}[>=stealth, thick, node distance=2.4cm]
    \node[state] (q) {};
    \node[draw, rectangle, minimum width=2.6cm, minimum height=1cm, right of=q, xshift=0.6cm] (N1) {};
    \node at ($(N1.west)+(0.5,0)$) {$q_1$};
    \node at ($(N1.east)+(-0.5,0)$) {$r_1$};
    \node[draw, rectangle, minimum width=2.6cm, minimum height=1cm, right of=N1, xshift=1.6cm] (N2) {};
    \node at ($(N2.west)+(0.5,0)$) {$q_2$};
    \node at ($(N2.east)+(-0.5,0)$) {$r_2$};
    \node[state, right of=N2, xshift=0.6cm] (r) {};

    \draw (q) edge[above] node{$\varepsilon$} (N1.west)
          (N1.east) edge[above] node{$\varepsilon$} (N2.west)
          (N2.east) edge[above] node{$\varepsilon$} (r);
\end{tikzpicture}
\end{center}
\noindent NFA for $L_1 \cdot L_2$.

\medskip
\noindent\textbf{For $L_1^*$:}
\begin{center}
\begin{tikzpicture}[>=stealth, thick, node distance=2.6cm]
    \node[state] (q) {};
    \node[draw, rectangle, minimum width=2.6cm, minimum height=1cm, right of=q, xshift=0.6cm] (N1) {};
    \node at ($(N1.west)+(0.5,0)$) {$q_1$};
    \node at ($(N1.east)+(-0.5,0)$) {$r_1$};
    \node[state, right of=N1, xshift=0.6cm] (r) {};

    \draw (q) edge[below] node{$\varepsilon$} (N1.west)
          (N1.east) edge[below] node{$\varepsilon$} (r)
          (N1.east) edge[bend left=60, above] node{$\varepsilon$} (N1.west);
\end{tikzpicture}
\end{center}
\noindent NFA for $L_1^*$. \hfill$\blacksquare$
\end{proof}

\section{Regular Expressions}

\begin{itemize}
    \item Algebraic representation of languages.
    \item Applications in text editors to find strings, compiler design.
\end{itemize}

\noindent Egs.\ of regular expressions:
\begin{center}
\begin{tabular}{c|l}
Regular Exp$^n$ & Languages \\ \hline
$0$ & $\{0\}$ \\
$\varepsilon$ & $\{\varepsilon\}$ \\
$0 \cup 01$ & $\{0, 01\}$ \\
$1^*$ & $\{\varepsilon, 1, 11, \dots\}$ \\
$(0 \cup 01) \cdot 1^*$ & $\{0, 01, 011, 0111, \dots\}$
\end{tabular}
\end{center}

\begin{greybox}
\begin{definition}
$R$ is said to be a \emph{regular expression} (RE) if $R$ has one of the following forms:
\begin{center}
\begin{tabular}{c|c|l}
Reg.\ Exp. & Language & Comment \\ \hline
$\phi$ & $\{\}$ & \\
$\varepsilon$ & $\{\varepsilon\}$ & \\
$a$ & $\{a\}$ & $a \in \Sigma$ \\
$R_1 \cup R_2$ & $L(R_1) \cup L(R_2)$ & \\
$R_1 \cdot R_2$ & $L(R_1) \cdot L(R_2)$ & \\
$R_1^*$ & $L(R_1)^*$ & \\
$(R_1)$ & $L(R_1)$ &
\end{tabular}
\end{center}
\end{definition}
\end{greybox}

\begin{greybox}
\begin{remark}
For every regular expression $R$, there is a unique language $L(R)$ corresponding to it.
\end{remark}
\end{greybox}

\begin{itemize}
    \item We often replace $\cup$ with the $+$ symbol and drop the $\cdot$ symbol.
    \item Precedence of the symbols/operators (High to Low): $(\,), \ *,\ \cdot,\ \cup$.
    \item $\phi^* = \{\varepsilon\}$.
\end{itemize}

\noindent\textbf{More Examples} (RE $\to$ language):
\begin{center}
\begin{tabular}{c|l}
Reg.\ Exp$^n$ & Language \\ \hline
$01$ & $\{01\}$ \\
$01 + 1$ & $\{01, 1\}$ \\
$(01 + \varepsilon) \cdot 1$ & $\{011, 01\}$ \\
$(0 + 10)^* \cdot (\varepsilon + 1)$ & $\{\varepsilon, 0, 1, 00, 01, 10, 000, 001, 101, 0000, \dots\}$
\end{tabular}
\end{center}

\begin{example}
(Language $\to$ RE)
\begin{enumerate}[(i)]
    \item $\{w \mid w \text{ has a single } 1\}$: \quad $R = 0^* \cdot 1 \cdot 0^*$.
    \item $\{w \mid w \text{ has at most one } 1\}$: \quad $R = 0^* \{\varepsilon, 1\} 0^* = 0^* + 0^*10^*$.
    \item $\{w \mid |w| \equiv 0 \pmod 3\}$:
    \[
    R = \{000, 001, \dots, 110, 111\}^* = ((0+1)(0+1)(0+1))^*.
    \]
    \item $\{w \mid w \text{ has a } 1 \text{ at every odd position and } |w| \text{ is odd}\}$:
    \[
    R = \{10, 11\}^* \{1\} = (1(0+1))^* \cdot 1.
    \]
\end{enumerate}
\end{example}

\subsection{Properties of RE}

\begin{enumerate}[1)]
    \item[1)] (a) $R_1 + (R_2 + R_3) = (R_1 + R_2) + R_3$ \\
    (b) $R_1(R_2 R_3) = (R_1 R_2) R_3$
    \item[2)] (a) $(R_1 + R_2) \cdot R_3 = R_1 \cdot R_3 + R_2 \cdot R_3$ \\
    (b) $R_1 \cdot (R_2 + R_3) = R_1 R_2 + R_1 R_3$
    \item[3)] $R_1 + R_2 = R_2 + R_1$ \quad (only $+$ is commutative)
    \item[4)] (a) $R_1 + \phi = R_1$ \\
    (b) $R_1 \varepsilon = R_1$
    \item[5)] $(R_1^*)^* = R_1^*$.
\end{enumerate}

\begin{blackbox}
\begin{theorem}
A language $L$ is regular iff there is a regular exp$^n$ $R$ s.t.\ $L = L(R)$.
\end{theorem}
\end{blackbox}

\begin{proof}[Proof (if part)]
This proof is based on case analysis.

\begin{center}
\begin{tabular}{c|c}
RE & NFA \\ \hline
$\phi$ & single non-accepting state \\
$\varepsilon$ & single start = accept state \\
$a$ & two states, start $\xrightarrow{a}$ accept
\end{tabular}
\end{center}

\noindent\textbf{$R_1 + R_2$:} Suppose NFA of $R_1$ \& $R_2$ resp.\ are
\begin{center}
\begin{tikzpicture}[thick]
    \node[draw, rectangle, minimum width=2.6cm, minimum height=1cm] (N1) {};
    \node at ($(N1.west)+(0.5,0)$) {$q_1$};
    \node at ($(N1.east)+(-0.5,0)$) {$r_1$};
    \node[below=0.2cm of N1] {$N_1$ for $L(R_1)$};

    \node[draw, rectangle, minimum width=2.6cm, minimum height=1cm, right=2cm of N1] (N2) {};
    \node at ($(N2.west)+(0.5,0)$) {$q_2$};
    \node at ($(N2.east)+(-0.5,0)$) {$r_2$};
    \node[below=0.2cm of N2] {$N_2$ for $L(R_2)$};
\end{tikzpicture}
\end{center}

\noindent $\therefore$ NFA for $L(R_1) \cup L(R_2)$ or $R_1 + R_2$ is
\begin{center}
\begin{tikzpicture}[>=stealth, thick, node distance=2.6cm]
    \node[state] (q) {};
    \node[draw, rectangle, minimum width=2.6cm, minimum height=1cm, above right of=q, xshift=1.2cm, yshift=0.3cm] (N1) {};
    \node at ($(N1.west)+(0.5,0)$) {$q_1$};
    \node at ($(N1.east)+(-0.5,0)$) {$r_1$};
    \node[draw, rectangle, minimum width=2.6cm, minimum height=1cm, below right of=q, xshift=1.2cm, yshift=-0.3cm] (N2) {};
    \node at ($(N2.west)+(0.5,0)$) {$q_2$};
    \node at ($(N2.east)+(-0.5,0)$) {$r_2$};
    \node[state, accepting, right of=N1, xshift=1.6cm, yshift=-1.1cm] (r) {};

    \draw (q) edge[above] node{$\varepsilon$} (N1.west)
          (q) edge[below] node{$\varepsilon$} (N2.west)
          (N1.east) edge[above] node{$\varepsilon$} (r)
          (N2.east) edge[below] node{$\varepsilon$} (r);
\end{tikzpicture}
\end{center}

\noindent\textbf{$R_1 R_2$:}
\begin{center}
\begin{tikzpicture}[>=stealth, thick, node distance=2.4cm]
    \node[state, initial] (q) {};
    \node[draw, rectangle, minimum width=2.6cm, minimum height=1cm, right of=q, xshift=0.6cm] (N1) {};
    \node[below=0.2cm of N1] {NFA for $L(R_1)$};
    \node[draw, rectangle, minimum width=2.6cm, minimum height=1cm, right of=N1, xshift=1.6cm] (N2) {};
    \node[below=0.2cm of N2] {NFA for $L(R_2)$};
    \node[state, accepting, right of=N2, xshift=0.6cm] (r) {};

    \draw (q) edge[above] node{$\varepsilon$} (N1.west)
          (N1.east) edge[above] node{$\varepsilon$} (N2.west)
          (N2.east) edge[above] node{$\varepsilon$} (r);
\end{tikzpicture}
\end{center}

\noindent\textbf{$R_1^*$:}
\begin{center}
\begin{tikzpicture}[>=stealth, thick, node distance=2.6cm]
    \node[state] (q0) {};
    \node[draw, rectangle, minimum width=2.6cm, minimum height=1cm, right of=q0, xshift=0.6cm] (N1) {};
    \node at ($(N1.west)+(0.5,0)$) {$q$};
    \node at ($(N1.east)+(-0.5,0)$) {$r$};
    \node[below=0.2cm of N1] {NFA for $L(R_1)$};
    \node[state, accepting, right of=N1, xshift=0.6cm] (r0) {};

    \draw (q0) edge[loop above] node{$\varepsilon$} (q0)
          (q0) edge[below] node{$\varepsilon$} (N1.west)
          (N1.east) edge[below] node{$\varepsilon$} (r0)
          (N1.east) edge[bend left=60, above] node{$\varepsilon$} (N1.west);
\end{tikzpicture}
\end{center}

So given an RE, we can explicitly construct the NFA.
\end{proof}

\section{Converting a DFA to a Regular Expression}

A GNFA is an NFA with the transitions being labelled by regular expressions.

\begin{greybox}
\begin{definition}
A GNFA is said to \emph{accept} a string $w$ if $w$ can be written as $w = w_1 \cdots w_k$ (concatenation of $k$ strings) and $\exists$ states $q_0, q_1, \dots, q_k$ in the GNFA such that
\begin{enumerate}[i)]
    \item $q_0$ is the start state.
    \item $q_k$ is an accept state.
    \item $\forall i$, if $R_i$ is the label on the transition $q_{i-1}$ to $q_i$, then $w_i \in L(R_i)$.
\end{enumerate}
\end{definition}
\end{greybox}

Without loss of generality, we may assume the following:
\begin{enumerate}[1.]
    \item GNFA has a unique accept state.
    \item There are no incoming transitions to the start state and no outgoing transitions from the accept state.
    \item There are transitions from the start state to every other state, and from every state to the unique accept state.
    \item There is a transition between every pair of states that are not the start/accept state.
\end{enumerate}

\begin{example}
Here is an example of a generalized NFA (GNFA), i.e.\ an NFA with regular-expression-labelled transitions.

\begin{center}
\begin{tikzpicture}[>=stealth, thick, node distance=3.5cm]
    \node[state, initial] (qs) {$q_S$};
    \node[state, right of=qs] (q1) {$q_1$};
    \node[state, below of=q1, yshift=1cm] (q2) {$q_2$};
    \node[state, accepting, right of=q1] (qacc) {$q_{acc}$};

    \draw (qs) edge[above] node{$01^*$} (q1)
          (q1) edge[loop above] node{$11$} (q1)
          (q1) edge[above] node{$0+01$} (qacc)
          (q1) edge[bend right, right] node{$(11)^*$} (q2)
          (q2) edge[bend right, left] node{$1$} (q1)
          (qs) edge[below] node{$\phi$} (q2)
          (q2) edge[loop below] node{$0$} (q2)
          (qs) edge[bend right=40, below] node{$00^*$} (qacc);
\end{tikzpicture}
\end{center}
\end{example}

\subsection{A Motivating Example}

\begin{center}
\begin{tikzpicture}[>=stealth, thick, node distance=3cm]
    \node[state, initial] (q1) {$q_1$};
    \node[state, right of=q1] (q2) {$q_2$};
    \node[state, accepting, below of=q2, yshift=1cm] (q3) {$q_3$};

    \draw (q1) edge[above] node{$a,b$} (q2)
          (q2) edge[loop above] node{$a$} (q2)
          (q1) edge[bend right, below] node{$a$} (q3)
          (q2) edge[bend left, right] node{$b$} (q3)
          (q3) edge[bend left, left] node{$b$} (q2);
\end{tikzpicture}
\end{center}

\noindent\textbf{Algorithm for DFA $\to$ RE:}
\begin{enumerate}[1.]
    \item Convert the DFA to a GNFA in the appropriate form.
    \item Removing a state.
\end{enumerate}

\medskip
\noindent\textbf{Step 1} (GNFA form):
\begin{center}
\begin{tikzpicture}[>=stealth, thick, node distance=3.2cm]
    \node[state, initial] (qs) {$q_S$};
    \node[state, right of=qs] (q1) {$q_1$};
    \node[state, right of=q1] (q2) {$q_2$};
    \node[state, accepting, right of=q2] (qa) {$q_A$};
    \node[state, below of=q2, yshift=1cm] (q3) {$q_3$};

    \draw (qs) edge[above] node{$\varepsilon$} (q1)
          (q1) edge[above] node{$a,b$} (q2)
          (q2) edge[loop above] node{$a$} (q2)
          (q2) edge[above] node{$\varepsilon$} (qa)
          (q1) edge[bend right, left] node{$a$} (q3)
          (q3) edge[bend right, right] node{$b$} (q2)
          (q2) edge[bend left, right] node{$b$} (q3)
          (q3) edge[bend right=50, above] node{$\varepsilon$} (qa)
          (qs) edge[bend left=20, above] node{$\phi$} (q2)
          (qs) edge[bend right=45, below] node{$\phi$} (qa)
          (qs) edge[bend right=60, below] node{$\phi$} (q3)
          (q1) edge[bend left=20, above] node{$\phi$} (qa);
\end{tikzpicture}
\end{center}

\medskip
\noindent\textbf{Step 2} (removing a state -- general picture): removing $q_{\text{rem}}$ from $q_A \xrightarrow{R_0} q_B$, $q_A \xrightarrow{R_1} q_{\text{rem}}$, $q_{\text{rem}} \xrightarrow{R_2} q_{\text{rem}}$ (self-loop), $q_{\text{rem}} \xrightarrow{R_3} q_B$ gives
\[
q_A \xrightarrow{\ R_0 + R_1 R_2^* R_3\ } q_B.
\]

\medskip
\noindent\textbf{Removing $q_1$:} the states $q_S, q_2, q_A, q_3$ remain, with (among others) the transition
\[
q_S \xrightarrow{\ \phi + \varepsilon \cdot \phi^*(a+b)\ } q_2, \qquad
q_S \xrightarrow{\ \phi + \varepsilon \cdot \phi^* \cdot \varepsilon\ } q_A, \qquad
q_S \xrightarrow{\ \phi + \varepsilon \cdot \phi^* \cdot \phi\ } q_3, \quad \text{etc.}
\]

\medskip
\noindent\textbf{Removing $q_2$:} with $q_S, q_A, q_3$ remaining,
\[
q_S \xrightarrow{\ \phi + (a+b) \cdot a^* \cdot \phi\ } q_3 = (a+b)a^*b, \qquad
q_3 \xrightarrow{\ \phi + (b + a(a+b)) \cdot a^* \cdot b\ } q_3 = (b + a(a+b))a^*b,
\]
\[
q_3 \xrightarrow{\ (\varepsilon + a) + (b + a(a+b))a^* \cdot \phi\ } q_A = \varepsilon + a.
\]

\medskip
\noindent\textbf{Removing $q_3$:} with $q_S, q_A$ remaining,
\[
q_S \xrightarrow{\ \varepsilon + (a+b)a^*b \cdot \big((b+a(a+b))a^*b\big)^* (\varepsilon + a)\ } q_A.
\]

\noindent Required RE:
\[
R = \varepsilon + (a+b)a^*b\big((b+a(a+b))a^*b\big)^*(\varepsilon + a).
\]

\section{More Closure Properties of Regular Languages}

\subsection*{Complement}

Complement of $L$: \quad $\overline{L} = \Sigma^* \setminus L$.

\begin{blackbox}
\begin{proposition}
If $L$ is regular, then $\overline{L}$ is also regular.
\end{proposition}
\end{blackbox}

\begin{proof}
Let $D = (Q, \Sigma, \delta, q_0, F)$ be some DFA for $L$. We construct a DFA $D'$ for $\overline{L}$ where
\[
D' = (Q, \Sigma, \delta, q_0, Q \setminus F).
\]
\end{proof}

\subsection*{Intersection}

\begin{blackbox}
\begin{proposition}
If $L_1$ and $L_2$ are regular, then $L_1 \cap L_2$ is also regular.
\end{proposition}
\end{blackbox}

\begin{proof}
From De Morgan's law, $L_1 \cap L_2 = \overline{\overline{L_1} \cup \overline{L_2}}$.

Now $L_1, L_2$ are regular $\implies \overline{L_1} \cup \overline{L_2}$ is regular $\implies L_1 \cap L_2$ is regular.
\end{proof}

\subsection*{Set Difference}

\begin{blackbox}
\begin{proposition}
If $L_1$ and $L_2$ are regular, then $L_1 \setminus L_2$ is also regular.
\end{proposition}
\end{blackbox}

\begin{proof}
$L_1, L_2$ regular $\implies \overline{L_1}, \overline{L_2}$ regular $\implies L_1 \cap \overline{L_2}$ regular $\implies L_1 \setminus L_2$ regular.
\end{proof}

\section{Reversal of a Language}

Let $w = a_1 a_2 \cdots a_n$ be a string, then
\[
\mathrm{rev}(w) = a_n a_{n-1} \cdots a_2 a_1.
\]

For a language $L \subseteq \Sigma^*$,
\[
\mathrm{rev}(L) := \{ w \in \Sigma^* \mid \mathrm{rev}(w) \in L \}.
\]

\begin{blackbox}
\begin{proposition}
If $L$ is regular, then $\mathrm{rev}(L)$ is also regular.
\end{proposition}
\end{blackbox}

\begin{proof}
Let $D = (Q, \Sigma, \delta, q_0, F)$ be a DFA for $L$. Now define an NFA $D'$ as the following.
\[
D' = (Q', \Sigma, \delta', q_0', F'), \qquad \text{where } Q' = Q \cup \{q_0'\}
\]
\[
\delta'(q, a) = \{ r \mid \delta(r, a) = q \}, \qquad
\delta'(q_0', \varepsilon) = \{ f \mid f \in F \}, \qquad
F' = \{q_0\}.
\]

Now $L(D') = \mathrm{rev}(L)$.

$\therefore\ \mathrm{rev}(L)$ is also regular.
\end{proof}

\section{Homomorphisms}

\begin{greybox}
\begin{definition}[Homomorphism]
Let $\Sigma, \Gamma$ be two alphabets. A \emph{homomorphism} $h : \Sigma \to \Gamma^*$ is a map such that if $w = a_1 a_2 \cdots a_n \in \Sigma^*$, then
\[
h(w) = h(a_1) h(a_2) \cdots h(a_n).
\]
\end{definition}
\end{greybox}

Also, $h(L) := \{ h(w) \mid w \in L \}$ where $L \subseteq \Sigma^*$.

\begin{blackbox}
\begin{proposition}
If $L$ is regular, then $h(L)$ is regular.
\end{proposition}
\end{blackbox}

\section{Inverse Homomorphism}

Let $h : \Sigma \to \Gamma^*$ be a homomorphism. For $L \subseteq \Gamma^*$, define
\[
h^{-1}(L) = \{ w \in \Sigma^* \mid h(w) \in L \}.
\]

\begin{blackbox}
\begin{proposition}
If $L$ is regular, then $h^{-1}(L)$ is regular.
\end{proposition}
\end{blackbox}

\begin{proof}
Proof of the above two propositions are via the same method: as $L$ is regular, we get a regular expression for $L$, and applying $h$ to each symbol in that RE, we get an RE for $h(L)$, thus it is regular. Same for $h^{-1}(L)$.
\end{proof}

\chapter*{Non-Regular Languages}
\addcontentsline{toc}{section}{Non-Regular Languages}

There exist non-regular languages. Non-regular languages require us to maintain some count. The following theorem helps us to sometimes check if a language is non-regular.

\begin{blackbox}
\begin{lemma}[The Pumping Lemma for Regular Languages]
Let $L$ be a regular language. Then $\exists\ p \geq 0$ s.t.\ $\forall w \in L$ with $|w| \geq p$, $\exists$ a partition of $w$ as $w = xyz$ such that
\begin{enumerate}[(i)]
    \item $|xy| \leq p$
    \item $y \neq \varepsilon$
    \item $\forall k \geq 0$, $xy^k z \in L$.
\end{enumerate}
\end{lemma}
\end{blackbox}

\begin{blackbox}
\begin{corollary}[Contrapositive of the Pumping Lemma]
$\forall p \geq 0$, $\exists w \in L$ with $|w| \geq p$ such that $\forall$ partitions $w = xyz$ with $|xy| \leq p$, $y \neq \varepsilon$, and $\exists k \geq 0$ s.t.\ $xy^k z \notin L$, then $L$ is not regular.
\end{corollary}
\end{blackbox}

\begin{example}
$L_1 = \{0^n 1^n \mid n \geq 0\}$ is not regular.

Take any $p \geq 0$. Choose $w = 0^p 1^p$, also take any partition of $w = xyz$ with $|xy| \leq p$ and $y \neq \varepsilon$.

$\because |xy| \leq p \implies 1 \notin x, 1 \notin y$, i.e., $x, y$ are made up of $0$'s only.

Choose $k = 0$, then $xy^0 z = 0^l 1^p$ where $l = p - |y| \neq p$.

$\therefore xy^0 z \notin L_1 \implies L_1$ is not regular. \hfill$\blacksquare$
\end{example}

\begin{example}
$L_2 = \{a^l b^m c^n \mid l + m \leq n\}$ is not regular. \quad $[\Sigma = \{a,b,c\}]$

Take $p \geq 0$. Choose $w = a^p b^p c^{2p}$. Given a partition $w = xyz$ s.t.\ $|xy| \leq p$ and $y \neq \varepsilon$, we have $y = a^t$ for some $t > 0$.

Choose $k = 2$: $xy^2 z = a^{p+t} b^p c^{2p}$.

Now $p + t + p \not\leq 2p \implies xy^2 z \notin L_2$

$\implies L_2$ is not regular.
\end{example}

\begin{example}
$L_3 = \{ w \in \{0,1\}^* \mid \#0 \text{ in } w = \#1 \text{ in } w \}$ is not regular.

Notice that $L(0^*1^*)$ is regular. Suppose for the sake of contradiction, $L_3$ is regular. Then
\[
L(0^*1^*) \cap L_3 \text{ is regular}
\]
\[
\implies L_1 = \{0^n 1^n \mid n \geq 0\} \text{ is regular.} \implies\Leftarrow
\]

$\therefore L_3$ is not regular. \hfill$\blacksquare$
\end{example}

\begin{example}
$L_4 = \{0^p \mid p \in \mathbb{P}\}$ (the set of primes) is not regular. \quad $[\Sigma = \{0\}]$

Given $q \geq 0$, choose $w = 0^r$ where $r \in \mathbb{P}$, $r \geq 2q$. Now given a partition of $w = xyz$ where $|xy| \leq q$ \& $|y| = t > 0$, choose $i = r+1$, then
\[
xy^i z = 0^{r(1+t)} \notin L_4.
\]

$\therefore L_4$ is not regular. \hfill$\blacksquare$
\end{example}

\subsection*{Proof of the Pumping Lemma for Regular Languages}

\begin{proof}
Let $D = (Q, \Sigma, \delta, q_0, F)$ be a DFA for $L$. Set $p = |Q|$.

\begin{center}
\begin{tikzpicture}[>=stealth, thick, node distance=3cm]
    \node[draw, rectangle, minimum width=3cm, minimum height=1.2cm] (D) {};
    \node[state, initial, left=0.1cm of D.west, anchor=west, xshift=0.3cm] (q0) {$q_0$};
    \node[state, accepting, right=0.1cm of D.east, anchor=east, xshift=-0.3cm] (qt) {$q_t$};
    \node[below=0.1cm of D.south] {$D$};
\end{tikzpicture}
\end{center}

Given $w \in L$ such that $|w| \geq p$, $\exists$ a sequence of states $q_0, \dots, q_t$ such that $t \geq p$ and $q_t \in F$.

By the Pigeonhole Principle, $\exists\ m, n \geq 0$ such that $0 \leq m < n \leq t$ \& $q_m = q_n$.

\begin{center}
\begin{tikzpicture}[>=stealth, thick, node distance=2.5cm]
    \node[state, initial] (q0) {$q_0$};
    \node[state, right of=q0, xshift=1cm] (qm) {$q_m = q_n$};
    \node[state, accepting, right of=qm, xshift=1cm] (qt) {$q_t$};

    \draw (q0) edge[above] node{$x$} (qm)
          (qm) edge[loop above] node{$y$} (qm)
          (qm) edge[above] node{$z$} (qt);
\end{tikzpicture}
\end{center}

Set $x$ = the string from $q_0$ to $q_m$, and let $m$ be the first time there is such a repetition. Set $y$ = string on the loop (first loop), $z$ = remaining string.

Now $\forall i$, $xy^i z \in L$.

$(\ast)$ all the $|w| \geq p$, $|xy| \leq p$, $y \neq \varepsilon$ take care of themselves from $p = |Q|$, $y$ pumped, etc. \hfill$\blacksquare$
\end{proof}

\chapter*{DFA Minimization}
\addcontentsline{toc}{section}{DFA Minimization}

DFA minimization is an algorithm that produces a minimal DFA from a given DFA.

\subsection*{Algorithm}

Input: $D = (Q, \Sigma, \delta, q_0, F)$.

\begin{enumerate}[1.]
    \item Construct a table of all $\{p, q\}$ pairs where $p, q \in Q$.
    \item Mark a pair $\{p, q\}$ if $p \in F$ and $q \notin F$, or vice versa.
    \item Repeat the following until no more pairs can be marked: mark $\{p,q\}$ if $\{\delta(p,a), \delta(q,a)\}$ is marked for some $a \in \Sigma$.
    \item Two or more states are equivalent if they are not marked.
\end{enumerate}

\begin{example}
\begin{center}
\begin{tikzpicture}[>=stealth, thick, node distance=3cm]
    \node[state, initial] (1) {$1$};
    \node[state, accepting, above right of=1] (2) {$2$};
    \node[state, accepting, below right of=1] (3) {$3$};
    \node[state, right of=2] (4) {$4$};
    \node[state, right of=3] (5) {$5$};
    \node[state, accepting, right of=4, yshift=-1.5cm] (6) {$6$};

    \draw (1) edge[above] node{$a$} (2)
          (1) edge[below] node{$b$} (3)
          (2) edge[bend left, above] node{$a$} (4)
          (2) edge[bend right, below] node{$b$} (5)
          (3) edge[bend right, above] node{$b$} (4)
          (3) edge[bend left, below] node{$a$} (5)
          (4) edge[above] node{$a,b$} (6)
          (5) edge[below] node{$a,b$} (6)
          (6) edge[loop right] node{$a,b$} (6);
\end{tikzpicture}
\end{center}

\begin{center}
\renewcommand{\arraystretch}{1.3}
\begin{tabular}{c|c|c|c|c|c}
1 & & & & & \\ \hline
$\times$ & 2 & & & & \\ \hline
$\times$ & & 3 & & & \\ \hline
$\times$ & $\times$ & $\times$ & 4 & & \\ \hline
$\times$ & $\times$ & $\times$ & & 5 & \\ \hline
$\times$ & $\times$ & $\times$ & $\times$ & $\times$ & 6
\end{tabular}
\end{center}

\begin{center}
\begin{tikzpicture}[>=stealth, thick, node distance=3.2cm]
    \node[state, initial] (1) {$1$};
    \node[state, accepting, right of=1] (23) {$2,3$};
    \node[state, right of=23] (45) {$4,5$};
    \node[state, right of=45] (6) {$6$};

    \draw (1) edge[above] node{$a,b$} (23)
          (23) edge[above] node{$a,b$} (45)
          (45) edge[above] node{$a,b$} (6)
          (6) edge[loop above] node{$a,b$} (6);
\end{tikzpicture}
\end{center}

This is the minimized DFA for the given DFA above.
\end{example}

\begin{example}
\begin{center}
\begin{tikzpicture}[>=stealth, thick, node distance=3cm]
    \node[state, initial, accepting] (1) {$1$};
    \node[state, above right of=1] (6) {$6$};
    \node[state, right of=6] (3) {$3$};
    \node[state, accepting, right of=3, xshift=1.5cm] (2) {$2$};
    \node[state, right of=2, xshift=1.5cm] (5) {$5$};
    \node[state, below right of=1] (4) {$4$};
    \node[state, right of=4] (8) {$8$};
    \node[state, below right of=2] (7) {$7$};

    \draw (1) edge[bend left, above] node{$a$} (6)
          (6) edge[bend left, below] node{$b$} (1)
          (6) edge[above] node{$a$} (3)
          (1) edge[bend right, below] node{$b$} (4)
          (4) edge[bend right, above] node{$a$} (1)
          (4) edge[above] node{$b$} (8)
          (8) edge[bend right=15, below] node{$a$} (4)
          (3) edge[below] node{$b$} (8)
          (3) edge[above] node{$a$} (2)
          (8) edge[above] node{$b$} (2)
          (2) edge[bend left, above] node{$a$} (5)
          (5) edge[bend left, below] node{$b$} (2)
          (2) edge[bend right, below] node{$b$} (7)
          (7) edge[bend right, above] node{$a$} (2)
          (7) edge[above] node{$a$} (5)
          (5) edge[bend left=40, above] node{$b$} (1);
\end{tikzpicture}
\end{center}

\begin{center}
\renewcommand{\arraystretch}{1.3}
\begin{tabular}{c|c|c|c|c|c|c|c}
1 & & & & & & & \\ \hline
$\times$ & 2 & & & & & & \\ \hline
$\times$ & & 3 & & & & & \\ \hline
$\times$ & $\times$ & & & 4 & & & \\ \hline
$\times$ & $\times$ & & & & 5 & & \\ \hline
$\times$ & $\times$ & $\times$ & $\times$ & $\times$ & $\times$ & 6 & \\ \hline
$\times$ & $\times$ & $\times$ & $\times$ & $\times$ & $\times$ & & 7 \\ \hline
$\times$ & $\times$ & $\times$ & $\times$ & $\times$ & & & 8
\end{tabular}
\end{center}

\[
\{1 \approx 2\}, \quad \{3 \approx 4\}, \quad \{3 \approx 5\}, \quad \{6 \approx 7\}, \quad \{6 \approx 8\}, \quad \{7 \approx 8\}.
\]

\noindent So the required minimized DFA is
\begin{center}
\begin{tikzpicture}[>=stealth, thick, node distance=4cm]
    \node[state, initial, accepting] (12) {$1,2$};
    \node[state, right of=12] (345) {$3,4,5$};
    \node[state, right of=345] (678) {$6,7,8$};

    \draw (12) edge[bend left, above] node{$a$} (345)
          (345) edge[bend left, below] node{$b$} (12)
          (345) edge[bend left, above] node{$b$} (678)
          (678) edge[bend left, below] node{$a$} (345)
          (12) edge[bend right=50, below] node{$b$} (678)
          (678) edge[loop above] node{$a,b$} (678);
\end{tikzpicture}
\end{center}
\end{example}